\setcounter{subsection}{0}
\section{ANÁLISIS DEL PROBLEMA}

El Capítulo~III mostró qué se sabe del dominio y qué ofrece el mercado, y concluyó que ninguna
solución convierte la lista del colegio en un kit completo. Este capítulo corresponde a la Fase~1
(OE1), en la etapa de comprensión del negocio de CRISP-DM \citep{chapman_crisp-dm_2000}. Con las entrevistas, la
encuesta a los acudientes y el análisis de las listas, y siguiendo la forma de especificación de
IEEE 830 \citep{ieee_830_1998}, se identifican los actores y se escriben los requerimientos, las
restricciones y la especificación funcional. Lo obtenido aquí es la entrada de la Fase~2, que diseña la
arquitectura en el Capítulo~V.\par

% =====================================================================
\subsection{Requerimientos}

En esta sección se resumen los requerimientos del sistema y se detallan los más importantes.
Para ello, primero se presentan los actores que intervienen en el problema y el proceso que
se siguió para levantar sus necesidades, y luego se explica cómo se organizaron y
priorizaron los requerimientos resultantes.\par

\subsubsection{Actores del sistema}

La caracterización del dominio identificó cuatro actores con intereses distintos frente al
sistema. La Tabla~\ref{tab:actores} los presenta junto con su relación con Tukit y su nivel
de poder e interés, es decir, qué tanto pueden influir en el éxito del sistema y qué tanto
les afecta su resultado. El colegio ocupa la posición más estratégica, porque controla el
acceso a las familias mediante la lista oficial y, a la vez, tiene interés en que sus
acudientes compren con menos dificultad.\par

\begin{table}[H]
\centering\footnotesize
\setlength{\tabcolsep}{4pt}
\renewcommand{\arraystretch}{1.25}
\caption{Actores del sistema y su relación con Tukit.}
\label{tab:actores}
\begin{tabularx}{\textwidth}{@{}p{2.6cm}Xp{3.2cm}@{}}
\toprule
\textbf{Actor} & \textbf{Relación con el sistema} & \textbf{Poder e interés} \\
\midrule
Acudiente & Usuario final que recibe, ajusta y paga el kit. & Poder medio, interés alto \\
Colegio & Publica la lista oficial y es el canal de confianza hacia las familias. & Poder
alto, interés alto \\
Equipo de Tukit & Opera el catálogo, revisa las listas y valida los renglones dudosos. &
Poder alto, interés alto \\
Estudiante & Influye en la elección del kit, por ejemplo con los personajes que prefiere,
aunque no usa el sistema. & Poder bajo, interés alto \\
\bottomrule
\end{tabularx}
\end{table}

\subsubsection{Proceso de levantamiento}

Los requerimientos se obtuvieron principalmente a partir de una encuesta a acudientes del
Colegio Agustiniano Ciudad Salitre, aplicada de forma presencial en el Día de la Familia en el
colegio con autorización del rector. Esta modalidad permitió reunir en pocas horas a familias de todos los niveles,
desde preescolar hasta la educación media, y conversar en paralelo con docentes y directivos
sobre la forma en que la institución elabora sus listas.\par

El cuestionario constaba de dieciocho preguntas, diecisiete cerradas y una abierta, y su
tiempo de respuesta tuvo una mediana de 4,6 minutos. Su diseño siguió un enfoque que profundiza en qué intenta lograr la persona cuando compra, más
allá del producto en sí, e incluyó preguntas sobre el tiempo, los desplazamientos y las
dificultades que le exige la compra. Las
preguntas abordaron el perfil del hogar, el tiempo y los lugares que demanda la compra, la
principal dificultad que enfrenta la familia, los criterios con que elige los productos, su
disposición a pagar un sobreprecio, la satisfacción del estudiante con sus útiles y la
confianza que le daría el respaldo del colegio. Se recibieron 65 respuestas, de las cuales
se analizaron 60, ya que las cinco restantes no autorizaron el tratamiento de sus datos,
conforme a la Ley 1581 de 2012 \citep{congreso_ley1581_2012}.\par

Los hogares encuestados pertenecen en su mayoría al estrato 4 (68,3\,\%), el 78,3\,\% tiene
un solo hijo en el colegio y el 78,3\,\% compra por aplicaciones móviles al menos algunas
veces al mes, lo que indica que la población está familiarizada con la compra en línea. A
partir de las respuestas, cada hallazgo se tradujo en uno o varios requerimientos del
sistema, como resume la Tabla~\ref{tab:encuesta-requerimientos}.\par

\begin{table}[H]
\centering\footnotesize
\setlength{\tabcolsep}{4pt}
\renewcommand{\arraystretch}{1.25}
\caption{Hallazgos de la encuesta a acudientes y requerimientos que se derivan de ellos
($n=60$).}
\label{tab:encuesta-requerimientos}
\begin{tabularx}{\textwidth}{@{}XX@{}}
\toprule
\textbf{Hallazgo} & \textbf{Requerimiento que se deriva} \\
\midrule
El 65\,\% dedica cuatro horas o más a la compra de los útiles de un hijo, y el 68,3\,\%
compra en tres o más lugares. &
Entregar la lista completa resuelta en un solo lugar, en lugar de producto por producto. \\
El 76,7\,\% recibe la lista con más de un mes de anticipación. &
El problema no es la falta de tiempo para comprar, sino el esfuerzo que exige hacerlo, de
modo que el sistema debe reducir ese esfuerzo. \\
La dificultad más mencionada (43,3\,\%) es no encontrar el producto o encontrarlo agotado, y
el 23,3\,\% menciona los libros y textos escolares. &
Considerar la disponibilidad al armar los kits e informar con claridad los renglones que no
se venden en la plataforma, como los textos que distribuyen las editoriales. \\
El precio es el primer criterio para el 38,3\,\%, las marcas exactas que pide el colegio
para el 31,7\,\%, la calidad para el 18,3\,\% y los personajes que prefiere el hijo para el
11,7\,\%. &
Armar los kits equilibrando costo, fidelidad de marca, calidad y personajes, y ofrecer un
kit para cada prioridad. \\
El 70\,\% pagaría como máximo un 5\,\% más por recibir el kit armado, y el 30\,\% aceptaría
un sobreprecio de 10\,\% o más. &
Recomendar por omisión el kit de menor costo, y sugerir uno de mayor calidad solo a quien
declare que acepta pagar más. \\
El 90\,\% confiaría más en la plataforma si el colegio la recomienda (58,3\,\% ``mucho
más''). &
Operar con un modelo B2B2C, en el que el colegio publica la lista oficial y respalda el
servicio. \\
El 58,3\,\% usaría el servicio con probabilidad alta (4 o 5 en una escala de 1 a 5). &
Validar la propuesta de valor con un piloto en el mismo colegio. \\
\bottomrule
\end{tabularx}
\end{table}

En las conversaciones con docentes y directivos se trataron la forma en que el colegio
elabora y publica sus listas y los inconvenientes que observa cuando las familias compran
productos distintos a los solicitados. De ellas surgieron los requerimientos de la
institución, como publicar la lista oficial de cada grado, garantizar que las familias
compren sobre esa versión y no exigir marcas ni proveedores específicos.\par

Como complemento, se analizaron las listas reales de diecisiete colegios del sector, lo que
permitió conocer la variedad de formatos y de formas de nombrar los útiles y definir los
requerimientos de lectura y normalización. Con toda esta información se elaboró la
especificación de requerimientos siguiendo el estándar IEEE 830 y su sucesor, ISO/IEC/IEEE
29148 \citep{ieee_830_1998,iso_29148_2018}.\par
\subsubsection{Organización y priorización}

Los requerimientos se agrupan en tres categorías. Los requerimientos funcionales describen lo
que el sistema debe hacer, como leer una lista o generar los kits, y se organizan según el
componente del sistema que los atiende. Los requerimientos no funcionales describen cómo
debe comportarse el sistema, por ejemplo en tiempo de respuesta, seguridad o facilidad de
uso, y se organizan según los atributos de calidad de la norma ISO/IEC 25010
\citep{iso_25010_2011}. Por último, los requerimientos propios del aprendizaje automático
describen las condiciones que deben cumplir los modelos de inteligencia artificial, como la
calidad mínima de la lectura de listas o el comportamiento ante casos de baja certeza.\par

Cada requerimiento se priorizó con el método MoSCoW, que lo clasifica como indispensable
(\emph{Must}), importante pero no crítico (\emph{Should}), deseable (\emph{Could}) o fuera
del alcance de esta versión (\emph{Won't}). La especificación reúne 132 requerimientos, de los
cuales 67 son funcionales, 49 no funcionales y 16 propios del aprendizaje automático. La priorización se hizo sobre los 98 de la versión 3 del SRS; la
Tabla~\ref{tab:moscow} (Anexo~\ref{anx:tablas}) muestra su distribución por componente.\par

% [movida al Anexo: tab:moscow]

La proporción de requerimientos indispensables es alta, y responde a la naturaleza del
sistema. Una obligación legal no puede clasificarse como opcional, y tampoco pueden serlo los
requerimientos que definen qué ocurre cuando algo sale mal, por ejemplo cuando un pago falla
o cuando la pasarela de pagos envía dos veces la misma confirmación. Sin ellos, el sistema
podría cobrar dos veces o dejar un pedido a medias, lo que resulta inaceptable en una
plataforma que trata datos de menores y recibe dinero de las familias.\par

\subsubsection{Requerimientos más importantes}
\label{sec:req-importantes}

La Tabla~\ref{tab:req-resumen} reúne los trece requerimientos que guían el diseño, ocho no
funcionales, con la característica de la norma ISO/IEC 25010 a la que pertenecen
\citep{iso_25010_2011}, y cinco funcionales, ordenados según el recorrido de una lista, desde su carga
hasta la compra. Los no funcionales que más pesan son los de protección de los datos de las familias
y de los menores, seguridad de los pagos y tiempo de respuesta, porque de ellos depende la confianza
del acudiente. La descripción y la justificación de cada uno están en el
Anexo~\ref{anx:requerimientos}, y el conjunto completo, en el SRS.

\begin{table}[htbp]
\centering
\scriptsize
\caption{Requerimientos más importantes del sistema.}
\label{tab:req-resumen}
\begin{tabularx}{\textwidth}{@{}llX@{}}
\toprule
ID & Característica & Qué exige \\
\midrule
RNF01 & Seguridad & No guardar datos de tarjeta ni bancarios; el pago se captura en la pasarela \\
RNF02 & Seguridad & Guardar solo lo indispensable del menor, cifrado, y quitar su nombre antes del modelo externo (Ley 1581) \\
RNF03 & Eficiencia de desempeño & Kits en 20~s o menos en el 95\,\% de los casos tras normalizar; proceso completo en 5~min o menos \\
RNF04 & Usabilidad & Aplicación web progresiva desde el celular, sin instalación \\
RNF05 & Eficiencia de desempeño & 50 usuarios simultáneos y unos 1.000 pedidos en el día pico, con error menor al 5\,\% \\
RNF06 & Escalabilidad & Crecer a 1.000 usuarios simultáneos replicando servicios, sin rediseño \\
RNF07 & Fiabilidad & 90\,\% de disponibilidad de 12:00 a 14:00 y de 18:00 a 22:00 \\
RNF08 & Fiabilidad & La falla de un servicio no tumba a los demás; sin lectura, se ofrece la carga manual \\
\midrule
RF01 & Funcional & Elegir la lista oficial o cargarla en foto, PDF o texto, y leerla en renglones con cantidad \\
RF02 & Funcional & Asociar cada renglón con su concepto pese a sinónimos y errores, y enviar a revisión lo incierto \\
RF03 & Funcional & Hasta cuatro kits (Ahorrador, Premium, Fan y A tu Medida), con máximo tres opciones por renglón \\
RF04 & Funcional & Cambiar productos, elegir el kit, pagar en línea y pedir el retracto (Ley 1480) \\
RF05 & Funcional & El colegio publica la lista oficial y el equipo la revisa antes de llegar a las familias \\
\bottomrule
\end{tabularx}
\end{table}

\subsection{Restricciones}
\label{sec:restricciones}

La restricción de mayor peso sobre el diseño es legal, no técnica. La Ley 1581 de 2012
\citep{congreso_ley1581_2012} establece el régimen general de protección de datos
personales en Colombia, y el Decreto 1377 de 2013 \citep{presidencia_decreto1377_2013} lo
desarrolla con un régimen especial para los datos de niños, niñas y adolescentes. Una lista
escolar es siempre la lista de un menor, de modo que la restricción alcanza todo el flujo y
no un módulo aislado. El sistema trata los datos para una sola finalidad, la prestación del
servicio, y atiende los derechos de los artículos 14 y 15 de la ley mediante un caso de
soporte dentro de los plazos legales.

La propuesta inicial era
anonimizar el nombre del menor antes de guardar la lista y antes de invocar el modelo
externo. El sistema construido protege ese dato de otra forma, por minimización: no lo recoge. No pide el
nombre legal; el apodo con que el acudiente identifica a cada hijo es opcional y se cifra en
reposo, y la percepción extrae solo los ítems y el grado (PV-03). Es mejor porque un dato
que no se recoge no puede llegar a los registros de operación, a los conjuntos de evaluación
ni al modelo de percepción, y no depende de que un reconocedor de entidades lo detecte.
Queda un caso: el nombre que figure en el documento cargado viaja con la página hacia la API
externa, y su depuración (PE-06) está pendiente.

La segunda restricción proviene del dominio de pagos. El sistema no almacena ningún dato del
medio de pago del acudiente: ni número de tarjeta, ni fecha de expiración, ni código de
seguridad. Solo guarda el identificador de la transacción y su estado, porque la captura del
medio de pago ocurre en el dominio de Wompi, que el sistema usa en su entorno de pruebas
detrás de un puerto de pasarela con un adaptador simulado (ADR 0033). Esta restricción
mantiene a Tukit fuera del alcance de cumplimiento de los datos de tarjeta. La confirmación
no depende de la respuesta sincrónica de la pasarela: llega por el aviso firmado de Wompi y
por la consulta de la transacción cuando el acudiente regresa, y ambos caminos pasan por un
mismo procedimiento idempotente, porque el aviso puede llegar duplicado o fuera de orden, o
perderse.

Las restricciones técnicas se derivan del ADR 0013, un monorrepositorio con despliegue
independiente por servicio. El sistema construido tiene seis desplegables con código, cinco
servicios de dominio y una puerta de enlace, empaquetados en contenedores estándar para que
el entorno de evaluación sea el mismo que el de desarrollo. Hoy se levantan con Docker
Compose sobre un único PostgreSQL 16, con una base lógica por servicio, RabbitMQ y Redis; la
orquestación productiva se decidirá con el despliegue de finales de noviembre de 2026. Cada
evento se escribe en una bandeja de salida en la misma transacción que el cambio de estado, y
RabbitMQ lo lleva a la cola de cada consumidor (ADR 0046). Sin RabbitMQ, los consumidores leen
las bandejas, y el sistema funciona completo con una sola base de datos, a costa de concentrar
en ella la coordinación.

% =====================================================================
\subsection{Especificación Funcional}
\label{sec:especificacion-funcional}

El SRS v6.1 define cuatro actores. El acudiente (padre, madre o tutor) es el comprador final:
carga la lista, declara sus preferencias, elige un kit y paga, y solo ve sus propias
solicitudes, pedidos y casos. El operador de Tukit revisa las listas y los kits: atiende la
cola de revisión y valida las líneas, pero no publica listas oficiales; también opera la
logística de los pedidos. El administrador de
Tukit tiene los permisos del operador y, además, publica las listas oficiales (CO-01), consulta
las solicitudes de kits y gestiona el catálogo, la logística y los casos de soporte desde el
panel de administración. El niño o la niña no opera ninguna pantalla: influye en la compra,
sobre todo en el kit Fan, pero es sujeto de datos protegidos y no usuario. La entrega física
tampoco tiene un actor propio: la hacen repartidores propios, y el despacho exige el nombre y
el celular de quien entrega (UM-03).

El flujo funcional encadena cuatro etapas. La lista escolar, fotografiada o entregada como
documento, pasa primero por la \emph{percepción}, que la convierte en texto estructurado
con un modelo multimodal, y luego por la \emph{normalización}, que asigna cada renglón de
texto libre a un concepto del catálogo mediante recuperación híbrida. Los renglones que la
normalización no cierra con certeza suficiente pasan a la \emph{resolución}. La
\emph{optimización}, por último, elige sobre el catálogo los productos que mejor equilibran
costo, fidelidad de marca, disponibilidad y personaje, con pesos derivados de lo que el
acudiente declara por cada hijo: cuánto quiere gastar y, si quiere, el personaje favorito
(PS-01). El resultado son los kits que el acudiente pidió, de los cuales lleva uno al pago:
el Ahorrador y el Premium siempre, el Fan si indicó un personaje y «A tu Medida» si eligió el
punto medio de gasto o declaró un tope.

La resolución concentra la dificultad de la normalización. Cuando la confianza de una
asignación cae por debajo del umbral de PE-04, el ítem queda marcado como ambiguo y espera
una decisión humana en la cola de revisión del panel (AD-01), donde cada corrección queda
registrada con su autor y su marca de tiempo (AD-02).

El
diseño preveía que una persona revisara solo los ítems por debajo del umbral. En el sistema
construido, el equipo revisa una vez cada lista nueva completa antes de publicarla, con el
documento original y la lectura del sistema lado a lado, y con un plazo de 4 h entre la
carga y la aprobación (PE-04). Es mejor porque detecta también los errores leídos con
confianza alta, que un filtro por ítem deja pasar, y porque la lista revisada sirve después a
todos los acudientes del mismo colegio y grado (NR-07). La validación de este flujo con
usuarios está planificada y pendiente (Sección~\ref{sec:oe5}).

La magnitud del problema la fijó el levantamiento de listas escolares de diecisiete
colegios: 3.608 ítems expresados en 1.451 formas distintas de nombrar los artículos, que el
catálogo reduce a 99 conceptos canónicos. Pocos conceptos concentran casi todo el volumen.
Según las métricas del catálogo, que registran 3.610 renglones, 3.600 de ellos con concepto
asignado, 25 conceptos cubren el 80,97\,\% de estos y 38 cubren el 90\,\%. Los totales de las
distintas fuentes difieren levemente sin una explicación registrada (anexo CRISP-DM). El
resto es una cola larga de variaciones superficiales de los mismos conceptos. Esa
distribución explica por qué ninguna solución de coincidencia exacta de cadenas resuelve el
problema y por qué la resolución es una parte estructural del flujo y no una salvedad.
